\section{The dictionary}
\label{app:dictionary}

\subsection{The D'Hoker--Kraus brane}
\label{app:dictionary-brane}

The background of section~\ref{sec:2.2} is the
uncharged magnetic brane of~\cite{DHoker:2009mmn,DHoker:2009ixq}. Their
action is $\frac{1}{16\pi G_5}\int\sqrt{-g}\,(R + 12 - F^2)$, so that their
Maxwell term corresponds to $\alpha = 2$ in our normalisation and their
Newton constant to $2\kappa^2 = 16\pi G_5$; their field $B$ is ours (the
$\mathcal{B} = \sqrt3 B$ of their boundary flux is a counting convention and
their Chern--Simons term vanishes identically for a constant field). We fixed
each factor by two independent routes, the Maxwell coefficient from their
Einstein equations and horizon constraints and the Newton constant from the
Schwarzschild free energy at $B = 0$, without reference to the observables
being compared. With this dictionary our fixed-temperature entropy shift
reproduces theirs, and their throat value $e^{2V} = B/\sqrt3$ is
\eqref{eq:throat} at $\alpha = 2$.

\subsection{The central charges}
\label{app:dictionary-charges}

In the conventions
of~\cite{Osborn:1993cr},
\begin{equation}
\langle T_{\mu\nu}(x)T_{\sigma\rho}(0)\rangle = \frac{C_T\,\mathcal{I}_{\mu\nu,\sigma\rho}}{x^{2d}}~, \qquad
\langle J_\mu(x)J_\nu(0)\rangle = \frac{C_J\,I_{\mu\nu}}{x^{2(d-1)}}~,
\end{equation}
with the current normalised by
the algebra of its charges, $[Q^a, Q^b] = i\epsilon^{abc}Q^c$, so that a doublet carries
$T^3 = \pm\tfrac12$. This is the boundary shadow of the bulk convention
$F^a = dA^a + \tfrac12\epsilon^{abc}A^b\wedge A^c$ of section~\ref{sec:2.1}. The two holographic normalisations quoted in
section~\ref{sec:4.4} are equation (3.15) of~\cite{Buchel:2009sk} at zero
Gauss--Bonnet coupling, with the action normalised as in their (2.1) and
their Planck length $\ell_P^{d-1} = \kappa^2$ and equations (30) and (54)
of~\cite{Freedman:1998tz}; the ratio $C_J/C_T = \tfrac{(d-1)(d-2)}{2d(d+1)}\alpha$
follows by division, and the map $\alpha \to \alpha\,(d-1)(d-2)/2d(d+1)$
turns $\alpha_\star(d) = 16/d(d-1)(d-2)$ into $8/d^2(d+1)$.

The check of section~\ref{sec:4.4} is the following. Romans' $SU(2)\times U(1)$
gauged supergravity~\cite{Romans:1985ps} is a consistent truncation around
the supergravity dual of any four-dimensional $\mathcal{N} = 2$
superconformal theory~\cite{Gauntlett:2007sm,Malek:2017njj}, and its $SU(2)$ gauge field
is dual to the $SU(2)_R$ current, whose two-point function is fixed by $c$
alone. In any four-dimensional conformal field theory
$C_T = 40c/\pi^4$~\cite{Osborn:1993cr},
and the $\mathcal{N} = 1$ superconformal Ward identity gives, for the
$\mathcal{N} = 1$ R-current, $C_R = 4c/\pi^4$, the supercurrent coefficient of
equations (3.43) and (11.12) of~\cite{Osborn:1998qu}, whose lowest component is the
R-current~\cite{Anselmi:1997am}. We decompose the $\mathcal{N} = 2$
R-symmetry as $R_{\mathcal{N}=1} = \tfrac13 r + \tfrac43 I_3$. The
combination $r - 2I_3$, under which one supercharge is neutral, sits in a
linear multiplet and is therefore orthogonal to the R-current. It follows
that $C_r = 8\,C_{I_3}$ for the $U(1)_r$ and $SU(2)_R$
Cartan currents, and hence
\begin{equation}
C_J(SU(2)_R) = C_{I_3} = \frac{3c}{2\pi^4}~, \qquad \frac{C_J(SU(2)_R)}{C_T} = \frac{3}{80}~,
\end{equation}
in every $\mathcal{N} = 2$ superconformal theory.

On the gravity side, the coefficient of the $SU(2)$ kinetic term in Romans'
Lagrangian, in the form of~\cite{Lu:1999bw} with the scalar at its
extremum $\mathcal{X} = 1$ and the gauge field rescaled to unit structure constant, is
$\alpha_{\rm SUSY} = L^2/4$; we obtained the same value from the
eleven-dimensional reduction of~\cite{Gauntlett:2007sm} and from Romans'
original conventions via~\cite{Ohl:2010au}. Then
$(3/20)(1/4) = 3/80$, which fixes both normalisations in $d = 4$ from the
field theory side. The scalar $\mathcal{X}$ is sourced by the $SU(2)$ field strength
at the same order as the metric, $(\Box + 4)\,\delta\mathcal{X} = -\tfrac1{24}F^aF^a$
at the linearised level in our normalisation, so $\mathcal{X}$ has
$m^2 = -4$ and $\Delta = 2$, which is why Einstein--$SU(2)$ theory with $SU(2)$
flux is not a truncation of Romans' theory.

The $d = 3$ check is the same construction one dimension down. On
the gravity side, minimal $\mathcal{N} = 2$ gauged supergravity in four
dimensions has the geometrised Maxwell normalisation $-F^2$ relative to
$R$ and a gravitino of charge $1/\ell$ on that field~\cite{Caldarelli:1998hg};
the field of unit supercurrent charge is $A/\ell$, so
$\alpha_{\rm SUSY} = 2\ell^2 = 2$. The reduction of eleven-dimensional
supergravity on any Sasaki--Einstein seven-manifold~\cite{Gauntlett:2007ma}
gives $\alpha = 1/2$ for its gauge field, under which the gravitino has
charge $1/2$ (by matching the action to~\cite{Caldarelli:1998hg}), and rescaling to unit charge gives $2$ again. The charge
normalisation is fixed independently by requiring that the
Bogomol'nyi--Prasad--Sommerfield Reissner--Nordström solution saturate
$\Delta = R$, the chiral-primary bound of the boundary theory. On the field
theory side, the three-dimensional $\mathcal{N} = 2$ superconformal Ward
identity writes both the $R$-current and the stress-tensor two-point
functions in terms of a single coefficient~\cite{Closset:2012ru},
which in the conventions of~\cite{Osborn:1993cr} gives $C_R/C_T = 1/6$ for
every such theory, and a free chiral multiplet with $R$-charges
$(\tfrac12, -\tfrac12)$ reproduces it by explicit Wick contraction,
$C_R = 1/16\pi^2$ and $C_T = 3/8\pi^2$. The holographic formulas at
$d = 3$ give $C_J/C_T = \alpha/12$, and $2/12 = 1/6$.
